
% Large Language Models (LLMs) have shown great potential in solving complex tasks, but their effectiveness is limited when dynamically applied to different scenarios. Orchestrating LLMs through agentic workflows facilitates collaboration and specialization among them, enhancing their ability to handle complex tasks. However, crafting such workflows traditionally relies on human experts, which is time-consuming and lacks the flexibility needed for task adaptation. Recent studies have attempted to automate the search for effective workflows, but they fall short in representing the diversity and complexity of different workflow types.

% In this paper, we aim to automatically craft agentic workflows for orchestrating LLMs to solve complex tasks dynamically. We model the task-solving process as a sequence of LLM-invoking nodes connected by edges (e.g., code-represented control flows), collectively forming a vast search space of possible workflows. However, given that both the code and node sets are countably infinite, the resulting search space is vast, posing a significant challenge in identifying effective workflows.
% To address these challenges, we present \model, a Monte Carlo Tree Search-based framework that efficiently navigates this vast search space to discover agentic workflow with minimal human design effort.
% Specifically, \model \parasummary{jiayi todo} to systematically explore possible workflows and identify the most effective sequence of operations. Experimental results across six widely-used benchmarks demonstrate the effectiveness of \model, the workflows crafted by \model outperforming handcrafted state-of-the-art methods by an average of \textbf{5.8\%} and exceeding other automatic workflow generation methods by \textbf{19.7\%}.

% 1 + 2 说明遇到的问题
% 大语言模型表现出跨领域解决复杂任务的强大潜力。为了全面发挥其能力，LLM往往在一个包含详细的instruction与操作顺序的agentic workflow中被调用，而这一workflow的获得需要消耗大量的human effort，难以scalability and generalizability. 

% 近期的工作试图自动化的生成agentic workflow的整体或部分并对其进行优化，然而这些工作的初始需要一定的人工设置，无法做到完整且有效的 agentic workflow 的自动生成。

% 为了解决这一问题，我们首先model workflow为a sequence of LLM-invoking nodes connected by edges (e.g., graph, neural network, code)以表示all possible workflows，并把 automatic workflow optimization 建模为一个在无限搜索空间中搜索最优解的问题。

% We introduce AFLOW, an automated framework for effective workflow search. AFLOW 在一个具备一系列预先定义好的node： operator 的搜索空间中 utilizes Monte Carlo Tree Search for iterative optimization of nodes and edges with 树状结构经验 与 execute feedback 直到获取 effective workflow. 

% Large language models (LLMs) have demonstrated remarkable potential in solving complex tasks across diverse domains, typically harnessed through agentic workflows comprising detailed instructions and operational sequences. However, creating these workflows demands significant human effort, limiting scalability and generalizability. Recent research has attempted to automate the generation and optimization of agentic workflows, but these approaches still require initial manual setup and fall short of fully automated, effective workflow generation. To address this challenge, we reformulate workflow optimization as a search problem over code-represented workflows, where LLM-invoking nodes are connected by edges. We introduce \textbf{AFlow}, an automated framework that efficiently explores this space using Monte Carlo Tree Search, iteratively refining workflows through code modification, tree-structured experience, and execution feedback. Empirical evaluations across six benchmark datasets demonstrate AFlow's efficacy, yielding a 5.7\% average improvement over state-of-the-art baselines. Moreover, AFlow enables smaller models to outperform GPT-4o on specific tasks at 4.55\% of its inference cost in dollars. The code will be made available as open-source upon publication. 

Large language models (LLMs) have demonstrated remarkable potential in solving complex tasks across diverse domains, typically by employing agentic workflows that follow detailed instructions and operational sequences. However, constructing these workflows requires significant human effort, limiting scalability and generalizability. Recent research has sought to automate the generation and optimization of these workflows, but existing methods still rely on initial manual setup and fall short of achieving fully automated and effective workflow generation. To address this challenge, we reformulate workflow optimization as a search problem over code-represented workflows, where LLM-invoking nodes are connected by edges. We introduce \modelbold, an automated framework that efficiently explores this space using Monte Carlo Tree Search, iteratively refining workflows through code modification, tree-structured experience, and execution feedback. Empirical evaluations across six benchmark datasets demonstrate \model's efficacy, yielding a 5.7\% average improvement over state-of-the-art baselines. Furthermore, \model enables smaller models to outperform GPT-4o on specific tasks at 4.55\% of its inference cost in dollars. The code is available at \href{https://github.com/FoundationAgents/AFlow}{https://github.com/FoundationAgents/AFlow}.



% Empirical evaluations across six benchmark datasets demonstrate AFLOW's efficacy, yielding a 5.7\% average improvement over state-of-the-art baselines. Furthermore, AFLOW enables smaller models to surpass GPT-4 performance on specific tasks at 4.55\% of the computational cost.

% Large Language Models (LLMs) have shown great potential  in solving complex tasks, but their effectiveness is limited when dynamically applied to different scenarios. Orchestrating LLMs through agentic workflows 

% While manually crafted agentic workflows have shown promise, they lack scalability and generalizability. 

% We introduce AFLOW, an automated framework for optimal workflow synthesis. AFLOW defines a compact search space encompassing cognitive heuristics and utilizes Monte Carlo Tree Search for iterative optimization of prompts, nodes, and edges. Empirical evaluations across six benchmark datasets demonstrate AFLOW's efficacy, yielding a 5.7\% average improvement over state-of-the-art baselines. Furthermore, AFLOW enables smaller models to surpass GPT-4 performance on specific tasks at 4.55\% of the computational cost.

% \yuyu{add specific strategies, e.g., ``utilizes heuristic-guided exploration and adaptive node expansion''} 

%\bang{I think the logic flow of abstract can be improved: first point out the limit of LLM agents in dynamically solving complex tasks in different scenarios, and say why this happens. Then say that we propose SOPtimizer inspired by the SOP concept of solving real-world complex tasks. Next, summarize the key contributions of SOPimizer, and finally show the experimental improvement to showcase the significant improvement achieved.}

% 这里还有一个点可以强调，我们找出的SOP形成了一个Pareto Front. 

%simple yet comprehensive framework defined by node-to-workflow representations for acquiring SOP. Recent advancements in prompt and workflow optimization can be considered specific instances of this framework, highlighting the inherent trade-off between human design costs and optimization efficiency in the automated acquisition of agentic SOPs. Inspired by this, we introduce \model, an optimization method that minimizes human effort while improving efficiency compared to existing approaches. Experimental results across multiple datasets demonstrate the effectiveness of \model.


%
%Orchestrating multiple Large Language Models (LLMs) as using workflows enhances effectiveness in solving complex tasks by enabling specialization and collaboration among LLMs. 

%Standardized Operating Procedure (SOP) represents the optimal workflow for a given task in human society, providing efficiency and consistency while requiring substantial human effort to obtain. 
%Similarly, this applies to Large Language Models (LLMs), where optimizing workflows can be seen as a process of obtaining SOPs.
%\parasummary{How to express that obtaining agency SOPs also requires a lot of human effort, which leads to the importance of automated obtain of SOPs}
%We propose a simple yet comprehensive structure defined by node-to-workflow representations for acquiring SOPs. Recent works on prompt optimization and workflow optimization can be viewed as specific instances of this structure, highlighting that the automated acquisition of agentic SOPs involves a trade-off between human design costs and optimization efficiency. Inspired by this, we introduce SOPtimizer, an optimization method that reduces human effort while enhancing efficiency compared to previous approaches. Experimental results across multiple datasets demonstrate the effectiveness of SOPtimizer.

% 核心观点 
% 1. Professionally designed workflows enhance LLM capabilities and applications.
% 2. 某个任务下最优的WorkFlow是这个任务的SOP，获取SOP需要较
% Workflow的优化问题可以被认为是SOP的获取问题（怎么去引出这个表述是一个难点，我明早问一下承霖哥）
% 3. 我们为Agentic SOP的获取定义了一个全流程且完备的结构，过去的所有Prompt优化工作可以被认为是这个结构中的一个特例。
% 3.5-4. 选择结构的不同部分优化会带来HumanDesign Cost与优化效率的Trade off
% 4. 在这个结构下，我们提出了一种有效且表现好的优化方法，能够兼顾Human Design（无需设计即可针对任何任务）Cost与优化效率（25步以内）
% 5. 通过在六个数据集上的实验，我们的优化方法表现 Blablabla

% Standardized Operating Procedures (SOPs) are essential for problem-solving but costly to create. 
% \yuyu{What is the relationship between SOP and LLM/Agent?
% Does the creation cost refer to the labor cost or the machine cost?}
% Therefore, we introduce SOPtimizer, an automated framework that generates and optimizes SOPs for any task \yuyu{"any task", did the experiment confirm this claim?} \mc{suggest delete} by representing tasks through modular actions and combining them into workflows. SOPtimizer applies critical thinking principles and uses feedback-driven optimization to refine these workflows. 
% \yuyu{From the above sentences, it is not intuitively clear how to achieve performance and cost optimization.}
% Unlike previous methods, it focuses on both performance and cost optimization. Experimental results across various tasks \yuyu{how many tasks and benchmarks? e.g., 10 tasks?} show that SOPtimizer-generated workflows outperform manually crafted SOPs and other automated approaches, reaching up to 95\% of the performance of domain-specific state-of-the-art methods, while consistently appearing on the Pareto front. 
